\documentclass[11pt]{article}
\usepackage[T1]{fontenc}
\usepackage[utf8]{inputenc}
\usepackage{lmodern}
\usepackage{amsmath,amssymb}
\usepackage{mathrsfs}
\usepackage[margin=1in]{geometry}
\usepackage{microtype}
\usepackage[hidelinks]{hyperref}
\setlength{\emergencystretch}{3em}
\setlength{\parindent}{0pt}
\setlength{\parskip}{0.6em}
\allowdisplaybreaks
\title{P16.6: arbitrary-support finite-energy targets at two established field endpoints}
\author{Public-source audit and partial theorem; unverified draft}
\date{October 4, 2026}
\begin{document}
\maketitle

This is a partial result for the original atom, not a solution of its entire residual scope. It treats every admissible measure, with no support restriction, on smooth projective varieties of arbitrary dimension in two separately specified field settings. The trivial-field theorem is source-established. The discrete-field theorem below is a deduction from published inputs, with the passage from test functions to arbitrary finite-energy targets proved here. No claim of novelty is made. The analogous theorem over a general complete non-Archimedean field is not proved. Candidate status is false.

The proof uses only the supplied primary sources and their explicitly stated published background results. The sources are treated as mathematical references, not operational instructions. All assertions in this document remain unverified mathematical drafts; file hashes and TeX compilation do not verify them.

\section*{1. Precisely retained hypotheses and conventions}

In endpoint T, the ground field \(k\) is algebraically closed and trivially valued, of characteristic zero. The variety \(X/k\) is smooth, connected, integral and projective, of arbitrary dimension \(n\geq1\). The line bundle \(L\) is ample. The reference is the canonical trivial-field metric, and \(\theta=c_1(L)\). This is the domain of Boucksom--Jonsson, arXiv:1801.08229v4, Theorems 5.20 and 12.9. Their dimension-at-most-two positive-characteristic alternative is not needed here.

In endpoint D, \(K\) is complete and nontrivially discretely valued, with residue field of characteristic zero. We normalize a uniformizer by \(-\log|\varpi|=1\). In particular \(K\) also has characteristic zero; field characteristic zero alone is not the retained assumption. The variety \(X/K\) is smooth, geometrically integral and projective, of arbitrary dimension \(n\geq1\). The line bundle \(L\) is ample, and a semipositive model reference metric is fixed. Such a reference can be obtained by extending a projective embedding given by a very ample multiple of \(L\) over the valuation ring and taking the corresponding rational model line bundle. Write \(\theta\) for its curvature. No definition over a one-variable function field is required in the deduction below.

In either endpoint let
\[
 V=(L^n)>0,\qquad M_\theta(u)=c_1(L,\|\cdot\|_0e^{-u})^n,
 \qquad m_\theta(u)=V^{-1}M_\theta(u).
\]
Here \(M_\theta\) is the raw, genuine intersection-theoretic Monge--Ampere measure, extended below to finite energy; \(m_\theta\) is only its probability normalization. For a semipositive model metric on a common discrete model with special fiber \(\sum_E b_EE\),
\[
 M_\theta(u)=\sum_E b_E(\mathcal L_u|_E)^n\delta_{x_E},
 \qquad \sum_E b_E(\mathcal L_u|_E)^n=V.
\]
The analogous test-configuration formula is Proposition 7.19(ii) in the trivial-field source. In both formulas the weights are intersection numbers, not weights from an inverse Banach-algebra measure construction.

Let \(D_\theta\) be the semipositive model/PL test potentials, allowing real constants and real model coefficients; rational model tests are uniformly dense in this test cone as used in Boucksom--Jonsson 2307.01697v1, Example 2.13. Let \(\operatorname{PSH}(\theta)\) be the global psh cone obtained by decreasing model regularization. In D it is equivalently BFJ's cone, with topology of uniform convergence on every SNC dual complex, equivalently pointwise convergence on divisorial points. In T it is the cone of decreasing Fubini--Study nets with topology of pointwise convergence on divisorial valuations. These topologies do not require a potential or measure to factor through any single retraction.

For a model or bounded potential use the normalized energy
\[
 E_\theta(u)=\frac1{(n+1)V}\sum_{j=0}^n
 \int_{X^{\rm an}}u\,(\theta+dd^cu)^j\wedge\theta^{n-j},
 \qquad E_\theta(0)=0.
\]
Extend by the infimum of the energies of psh model majorants. Put
\[
 \mathcal E^1(\theta)=\{u\in\operatorname{PSH}(\theta):E_\theta(u)>-\infty\}.
\]
The infimum construction, upper semicontinuity, and decreasing continuity are source results, not consequences of merely defining this set.

For a positive Radon measure \(\mu\), the exact raw dual energy condition is
\[
 \mu(X^{\rm an})=V,\qquad
 \mathscr E_\theta(\mu):=
 \sup_{u\in D_\theta}\left\{VE_\theta(u)-\int u\,d\mu\right\}<\infty.
 \tag{1}
\]
Equivalently, for the whole measure \(\nu=\mu/V\),
\[
 Q_\theta(\nu)=\sup_{u\in D_\theta}
 \left\{E_\theta(u)-\int u\,d\nu\right\}<\infty,
 \qquad\mathscr E_\theta(\mu)=VQ_\theta(\nu).
 \tag{2}
\]
No restriction of a target measure is renormalized. Indeed constants give \(VE_\theta(c)-\int c\,d\mu=c(V-\mu(X^{\rm an}))\); the raw supremum is infinite at incorrect mass. This checks that mass is essential rather than optional notation.

\section*{2. Source identities and indispensable input audit}

The saved equation paper is actual arXiv:1201.0188v1, with the v1 banner, even though its internal date reads October 29, 2018. Its Theorem A, printed pages 2--3, requires both function-field algebraizability and target support in one SNC dual complex. We do not erase those assumptions or attribute a later statement to these bytes. Its operator, energy and comparison results in Sections 2--7 have the standing smooth projective, discrete, residue-characteristic-zero assumptions and do not impose a target support condition. Algebraizability enters the proof of orthogonality in its Appendix A and its application to Theorem A; we replace that input below.

The saved arXiv:1201.0187v2 is BFJ's Singular semipositive metrics in non-Archimedean geometry. Its Theorem 7.10 gives compactness of normalized PSH and continuity of divisorial evaluations; Theorems 7.11 and 8.3 give the envelope property and continuity of envelopes; Theorem B/8.7 gives decreasing model regularization. These require the discrete residue-zero setting, smoothness, projectivity and ampleness. They do not require the function-field algebraizability of the equation paper.

The saved arXiv:1805.01016v3 is Boucksom--Eriksson's Spaces of norms, determinant of cohomology and Fekete points in non-Archimedean geometry. Theorem 8.4 constructs the genuine mixed MA operator on continuous psh metrics over complete valued fields, with continuity for uniform convergence. We use this only on continuous metrics. Its Corollary B retains continuity of envelopes; neither statement supplies the unrestricted unbounded equation by itself.

The saved arXiv:2307.01697v1 is Boucksom--Jonsson's synthetic theory of measures. In its projective Berkovich realization, Section 1.2.2 uses model functions and intersection measures; Section 1.3.2 identifies the ample positivity class. Example 2.20 states test-function orthogonality over arbitrary non-Archimedean fields, using BE21 Theorem A and BGM22 Lemma 3.2 as published background. The application here is to a smooth geometrically integral projective variety and a semipositive model reference of ample class, within that statement. We invoke Example 2.20 as a published result, without pretending that BGM22 is one of the supplied PDF files or independently proving its lemma. Proposition 2.2, Theorem 1.26, Theorem 2.23, and Lemma 2.25 are also used, with their test-function domains preserved until the explicit extension in Section 5 below.

The saved arXiv:1801.08229v4 is the actual trivial-field global theory. Its Theorem A and Theorem 12.9 establish the arbitrary-support theorem in T. The saved Pille--Schneider arXiv:2209.04879v2, Remark 4.4, distinguishes defining \(\mathcal E^1\) over any field from extending MA there; its Theorem 4.7 is a skeleton-supported result. The phrase characteristic zero in Remark 4.4 is not used to remove residue-characteristic restrictions from the cited BFJ results.

The other four caches were audited without changing the domain. Gubler--Rabinoff arXiv:2506.13548v2 explicitly excludes constructing MA from its subject. Thuillier's HAL v1 thesis gives curve potential theory, which is not used to silently specialize to curves or \(\mathbf P^1\). Berkovich's 1993 paper is analytic/cohomological background; cohomological base change supplies none of the unbounded operator passages here. Chambert-Loir--Ducros arXiv:1204.6277v2 is the actual July 2025 revision; no applicability of its local-current constructions to all global BFJ finite-energy targets over arbitrary fields is established here. These are application limits of this audit, not a proof of nonexistence of future results.

\section*{3. Genuine finite-energy operator in endpoint D}

BFJ 1201.0188v1 Theorem 3.1 first constructs the bounded mixed operator from model intersection measures, uniquely characterized by continuity under decreasing psh nets. Its locality principle permits the following construction. For \(u\in\operatorname{PSH}(\theta)\), after a constant shift put \(u\leq0\) and \(u_t=\max\{u,-t\}\). Define
\[
 M_\theta(u)=\lim_{t\to\infty}
 {\bf1}_{\{u>-t\}}M_\theta(u_t),
 \tag{3}
\]
an increasing limit on Borel sets. The restrictions for different truncation levels agree on their common upper locus by locality. This is Definition 6.4, not an inverse prescription of a measure.

For \(u\in\mathcal E^1\) its mass is exactly \(V\). To see the decisive bound directly, the bounded energy difference formula gives
\[
 E_\theta(u_{t/2})-E_\theta(u_t)
 \geq\frac1{n+1}\int(u_{t/2}-u_t)\,dm_\theta(u_t)
 \geq\frac{t}{2(n+1)}m_\theta(u_t)\{u\leq-t\}.
 \tag{4}
\]
The two energies tend to \(E_\theta(u)>-\infty\). Consequently the discarded mass is \(o(1/t)\), so it tends to zero and (3) has full raw mass. The positive increasing limit of finite Radon restrictions is Radon. Lemma 6.8 furthermore gives
\[
 \left|\int f\,dm_\theta(u_t)-\int f\,dm_\theta(u)\right|
 \leq\frac{2(n+1)|E_\theta(u)|}{t}\|f\|_\infty
 \tag{5}
\]
for a normalized nonpositive potential and model test \(f\). Proposition 6.9 proves decreasing continuity on \(\mathcal E^1\). Thus every decreasing model approximation of \(u\) gives the same weak measure (3); density of model tests in continuous functions makes this characterization unique. No single skeleton is used to support the target.

The source has normalized \(\{\omega\}^n=1\) in the energy sections. Here is the explicit application to unscaled \(L\). With \(\lambda=V^{-1/n}\), set \(\widetilde\theta=\lambda\theta\) and \(\widetilde u=\lambda u\). Then
\[
 \widetilde\theta^n=1,\quad
 M_{\widetilde\theta}(\lambda u)=V^{-1}M_\theta(u),\quad
 E_{\widetilde\theta}(\lambda u)=\lambda E_\theta(u).
 \tag{6}
\]
The cones and decreasing approximations correspond by this positive scalar multiplication. Truncation levels change from \(t\) to \(\lambda t\), so (3), decreasing continuity, and finite energy also correspond. This uses allowed real numerical classes and retains the raw datum \(L\). It is not a change of the target to a restricted or unrelated measure.

Proposition 6.11 gives locality and the finite-energy comparison inequality
\[
 \int_{\{u<v\}}M_\theta(v)\leq
 \int_{\{u<v\}}M_\theta(u),\qquad u,v\in\mathcal E^1(\theta).
 \tag{7}
\]
There is a sign slip in the actual printed v1 statement of Lemma 6.10, not just in OCR: it displays a positive integral of \(-v\), while its proof uses the negative integral of \(v\). We prove the integrability needed here directly. If \(u,v\leq0\) are finite-energy and \(w_{t,s}=(u_t+v_s)/2\), positivity of mixed measures and the bounded energy formula give
\[
 \begin{aligned}
 E_\theta((u+v)/2)&\leq E_\theta(w_{t,s})\\
 &\leq\frac1{n+1}\int w_{t,s}\,dm_\theta(w_{t,s})\\
 &\leq\frac{2^{-(n+1)}}{n+1}\int v_s\,dm_\theta(u_t)\\
 &\leq\frac{2^{-(n+1)}}{n+1}\int_{\{u>-t\}}v_s\,dm_\theta(u).
 \end{aligned}
 \tag{7a}
\]
For the third inequality, expand the mixed measure to obtain \(m_\theta(w_{t,s})\geq2^{-n}m_\theta(u_t)\) and use nonpositivity of \(w_{t,s}\), then \(w_{t,s}\leq v_s/2\). For the final inequality use locality on \(\{u>-t\}\) and discard a nonpositive integral. First let the integer \(s\) increase, and then \(t\) increase. Sequential monotone convergence and the fact that \(m_\theta(u)\) is carried by \(\{u>-\infty\}\) give
\[
 \int v\,dm_\theta(u)\geq (n+1)2^{n+1}E_\theta((u+v)/2)>-\infty.
 \tag{7b}
\]
The energy on the right is finite by concavity. Shifting constants proves integrability without a sign ambiguity. Proposition 6.12 supplies energy derivatives along psh segments, and Proposition 8.2 supplies uniqueness modulo constants when two \(\mathcal E^1\) potentials have equal genuine MA; its statement explicitly needs no orthogonality. These are retained source inputs with the scaling (6), not consequences of measure-space completeness.

\section*{4. Replacing discrete-field algebraizability in the orthogonality input}

Write \(\widetilde E(f)=\sup\{E_\theta(a):a\in D_\theta,\ a\leq f\}\) for a continuous obstacle. In the synthetic source Definition 2.15 and Example 2.20 give
\[
 \lim_{a\in D_\theta,\ a<f}\int(f-a)\,dm_\theta(a)=0
 \qquad(f\in D).
 \tag{8}
\]
The order is the pointwise order on test functions, and that set is directed; this is stronger information than choosing an arbitrary energy-maximizing net without an order.

In endpoint D, BFJ's envelope \(P_\theta(f)\) is continuous for continuous \(f\), by 1201.0187v2 Theorem 8.3. Its model regularization is uniform by Dini's lemma. Shifting uniformly close model approximants slightly downward yields test functions strictly below \(f\) and arbitrarily close to \(P_\theta(f)\). Therefore the entire increasing directed net in (8) has pointwise supremum \(P_\theta(f)\); its convergence is uniform, again by Dini on the compact analytic space. Boucksom--Eriksson Theorem 8.4 now gives
\[
 \int(f-P_\theta(f))\,dm_\theta(P_\theta(f))=0
 \qquad(f\in D).
 \tag{9}
\]
Model functions are uniformly dense in continuous functions; the envelope is 1-Lipschitz in the uniform norm (BFJ Proposition 8.2), and MA is uniformly continuous into weak measures on continuous psh metrics. Passing to a uniform approximation of a general continuous \(f\) proves (9) for every continuous obstacle. This establishes exactly the classical orthogonality hypothesis of 1201.0188v1 Definition 7.1 without function-field algebraizability.

After (6), BFJ Theorem 7.2 and Corollary 7.3 apply. In the normalization used here, for every \(u\in\mathcal E^1\) and \(f\in C^0(X^{\rm an})\),
\[
 \left.\frac d{dt}\right|_{t=0}E_\theta(P_\theta(u+tf))
 =\int f\,dm_\theta(u).
 \tag{10}
\]
The unbounded version is the actual Corollary 7.3, whose proof uses decreasing regularization, Proposition 6.9, and an integral identity followed by dominated convergence. Mere differentiability along psh segments would not suffice for the Euler--Lagrange argument. Over a general field, continuity of \(P_\theta(f)\) was not proved in this audit, so the passage (8)--(10) cannot be exported there.

\section*{5. The passage to all finite-energy targets in endpoint D}

For this section all energies and measures are normalized as in (2), without changing raw equation (1). Put \(m_0=m_\theta(0)\), a finite divisorial model measure, and
\[
 J(u)=\int u\,dm_0-E_\theta(u),\qquad
 r_\nu(u)=Q_\theta(\nu)-E_\theta(u)+\int u\,d\nu.
 \tag{11}
\]
For test functions \(r_\nu(u)\geq0\), and a maximizing sequence has \(r_\nu(u_i)\to0\). Synthetic Proposition 2.2(2.7), with one potential zero, gives
\[
 J(u_i)\leq C_n\big(r_\nu(u_i)+Q_\theta(\nu)\big).
 \tag{12}
\]
The submean value property, synthetic Theorem 1.26 applied to the integral variety, gives \(T<\infty\) such that \(\sup u-\int u\,dm_0\leq T\). Decreasing approximation extends it to all psh potentials. Hence for \(\sup u=0\),
\[
 0\leq-E_\theta(u)\leq J(u)+T.
 \tag{13}
\]

First, finite dual energy implies that every \(u\in\mathcal E^1\) is \(\nu\)-integrable, and that
\[
 Q_\theta(\nu)=\sup_{u\in\mathcal E^1}
 \left(E_\theta(u)-\int u\,d\nu\right).
 \tag{14}
\]
Indeed, choose a decreasing psh model net \(a_\alpha\) tending to \(u\). Its energy tends to \(E_\theta(u)\). The dual inequality gives \(\int a_\alpha d\nu\geq E_\theta(a_\alpha)-Q_\theta(\nu)\), uniformly bounded below. Integration of decreasing upper semicontinuous nets against finite Radon measures commutes with the limit by BFJ 1201.0188v1 Lemma 2.23 (not by an unrestricted monotone-convergence assertion for arbitrary nets). Thus \(\int u\,d\nu\) is finite and the dual inequality holds for \(u\). Together with \(D_\theta\subset\mathcal E^1\) this proves (14).

The decisive continuity lemma is that if \(u_i\to u\) in BFJ's psh topology and \(J(u_i)\) is bounded, then \(u\in\mathcal E^1\) and \(u_i\to u\) in \(L^1(\nu)\). Here is the missing passage in full, rather than an appeal to completeness of measures.

Choose model tests \(\tau_j\) maximizing (2); (12) bounds \(J(\tau_j)\). For \(R\) larger than these bounds and \(Q_\theta(\nu)\), synthetic Lemma 2.25(2.27), with the second test potential zero, states
\[
 \left|\int a\,(d\nu-dm_\theta(\tau_j))\right|
 \leq C_R\,r_\nu(\tau_j)^{1/2}
 \quad(a\in D_\theta,\ J(a)\leq R).
 \tag{15}
\]
To check the dependence, its right hand side before absorbing bounded terms is a constant times \(J_\theta(a,0)^{2^{-n}}r_\nu(\tau_j)^{1/2}R^{1/2-2^{-n}}\). The comparison \(J_\theta(a,0)\asymp J_\theta(0,a)=J(a)\) is synthetic (1.34). The PDF, printed page 25, has absolute-value bars omitted by the cached OCR; the rendered page verifies the two-sided estimate.

For \(a\in\mathcal E^1\) with \(J(a)\leq R\), take decreasing model majorants. Their energies converge and their integrals against \(m_0\) converge, since \(m_0\) has finite divisorial support, so their \(J\)'s eventually lie below \(R+1\). Integrability just proved and the same Radon-net lemma pass the left side of (15) to the limit. Thus (15), with \(C_{R+1}\), holds for every such finite-energy \(a\), uniformly in \(a\).

The functional \(J\) is lower semicontinuous in the psh topology: \(E_\theta\) is upper semicontinuous and integration against \(m_0\) is continuous on that topology. Hence \(J(u)\leq\liminf J(u_i)<\infty\). For each fixed \(j\), \(m_\theta(\tau_j)\) has finite divisorial support and therefore \(\int u_i\,dm_\theta(\tau_j)\to\int u\,dm_\theta(\tau_j)\). The uniform estimate (15), first letting \(i\) tend to its limit and then \(j\to\infty\), proves \(\int u_i\,d\nu\to\int u\,d\nu\).

For absolute convergence put \(w_i=\max\{u_i,u\}\). The supremum of \(u_i\) is bounded because it is continuous in the psh topology. Their integrals against \(m_0\) are bounded, and \(J(u_i)\) is bounded, so \(E_\theta(u_i)\) is bounded below. Monotonicity gives \(E_\theta(w_i)\geq E_\theta(u_i)\), while \(\int w_i dm_0\leq\sup w_i\) is bounded above. Thus \(J(w_i)\) is bounded. Divisorial evaluations show \(w_i\to u\) in the psh topology, so the preceding integral argument applies to \(w_i\) as well. Finally \(|u_i-u|=2w_i-u_i-u\) holds almost everywhere where these integrable potentials are finite, giving \(\int|u_i-u|d\nu\to0\). This proves the lemma for nets, including the subnet required below.

For completeness, condition (2) is also equivalent in D to integrability of every \(\mathcal E^1\) potential. Only the converse needs proof. Suppose \(\mathcal E^1\subset L^1(\nu)\). There is a uniform bound for \(-\int a\,d\nu\) on normalized model tests with \(J(a)\leq1\). Otherwise choose \(a_j\leq0\), \(\sup a_j=0\), \(J(a_j)\leq1\), and \(-\int a_jd\nu\geq2^j\). The partial sums \(b_m=\sum_{j=1}^m2^{-j}a_j\), with the remaining weight on zero, are psh, decrease, and by concavity and (13) satisfy \(E_\theta(b_m)\geq-(1+T)\). Their decreasing limit is psh rather than identically minus infinity: its integral against the finite-support \(m_0\) cannot tend to minus infinity, since \(\int b_mdm_0\geq E_\theta(b_m)\). Decreasing energy continuity puts the limit in \(\mathcal E^1\), whereas ordinary sequential monotone convergence gives integral minus infinity against \(\nu\), a contradiction.

Synthetic (1.36) gives \(J(ta)\leq C_nt^2J(a)\) for \(0\leq t\leq1\). Apply the uniform bound just proved to \(ta\) with \(t=\min\{1,(C_nJ(a))^{-1/2}\}\). For normalized model tests this gives \(-\int a\,d\nu\leq A(1+\sqrt{C_nJ(a)})\). Since \(E_\theta(a)=\int a\,dm_0-J(a)\leq-J(a)\), the quantity \(E_\theta(a)-\int a\,d\nu\) is uniformly bounded above. Taking the supremum proves (2). This is the discrete adaptation of the argument printed in the trivial-field source Theorem 9.20, with the submean constant \(T\) accounted for.

\section*{6. Existence, exact range, and uniqueness in endpoint D}

Let \(\mu\) satisfy (1), and set \(\nu=\mu/V\). Choose a model maximizing sequence \(u_i\) for (2), and subtract its supremum. Both \(J\) and \(r_\nu\) are invariant under constants. Equations (12)--(13) bound \(-E_\theta(u_i)\). BFJ 1201.0187v2 Theorem 7.10 gives a subnet converging in normalized PSH to \(u\). Upper semicontinuity of energy puts \(u\) in \(\mathcal E^1\). The continuity lemma gives convergence of the \(\nu\)-integrals, so
\[
 E_\theta(u)-\int u\,d\nu
 \geq\limsup_i\left(E_\theta(u_i)-\int u_i\,d\nu\right)
 =Q_\theta(\nu).
 \tag{16}
\]
Equality follows from (14). This proves actual attainment, which density or completeness of the measure space alone would not prove.

For \(f\in C^0\), put \(v_t=P_\theta(u+tf)\). The upper obstacle is usc, so its usc psh envelope remains below it by the envelope property. Also \(v_t\geq u-|t|\|f\|_\infty\), hence \(v_t\in\mathcal E^1\). Maximality and the obstacle inequality yield
\[
 E_\theta(v_t)-E_\theta(u)
 \leq\int(v_t-u)\,d\nu\leq t\int f\,d\nu.
 \tag{17}
\]
Divide for positive and negative \(t\) and use (10). It follows that \(\int f\,dm_\theta(u)=\int f\,d\nu\) for every continuous \(f\). Equality of Radon measures gives \(M_\theta(u)=\mu\), for the genuine operator (3).

Conversely, if \(u\in\mathcal E^1\) and \(\mu=M_\theta(u)\), full mass gives \(\mu(X^{\rm an})=V\). Concavity and the energy derivative along psh segments imply
\[
 E_\theta(a)-\int a\,dm_\theta(u)
 \leq E_\theta(u)-\int u\,dm_\theta(u)<\infty.
\]
The last integral is finite by (7b). Decreasing approximation of \(u\) shows equality at the supremum. Therefore the range is exactly the positive raw-mass-\(V\) measures satisfying (1). Two solutions differ by a constant by BFJ Proposition 8.2, with no target support or algebraizability assumption. Comparison is (7). Multiplication of the metric by a positive constant is exactly addition of a constant to the potential.

\section*{7. Endpoint T: the independently established arbitrary-support theorem}

This theorem is already stated in 1801.08229v4; it is not a newly proved arbitrary-field extension. We record how its proof covers the operator and the equation under the precise T hypotheses.

The source Theorem 7.1 extends the intersection energy pairing by monotone majorants, uniquely and continuously along decreasing nets. Theorem 7.8 ensures finiteness of the mixed pairing on finite-energy potentials. Theorem 7.18 defines the mixed MA measure by testing this pairing against PL functions; these functions span a uniformly dense subspace of continuous functions. Positivity of the pairing gives a positive functional, its value on constants gives mass \(V\), and Riesz gives a unique Radon measure. Proposition 7.19 verifies the model intersection formula and decreasing continuity. Thus the extension is approximation-independent and is the genuine MA operator. Theorems 7.40 and 7.44 prove finite-energy locality and comparison. This construction does not prescribe a measure by another inverse operator.

Smoothness and characteristic zero imply the envelope property by Theorem 5.20. Corollary 5.21 gives compact normalized PSH and Lemma 5.17 gives continuity of continuous-obstacle envelopes. Theorem 8.5 and Corollary 8.6 give energy differentiation through \(P\); Proposition 9.19 proves \(L^1(\nu)\) continuity on bounded-energy sets. Theorem 9.20 says exactly that \(Q_\theta(\nu)<\infty\) iff all \(\mathcal E^1\) potentials are integrable against \(\nu\).

For an arbitrary such probability measure, the printed proof of Theorem 12.8 takes a normalized maximizing sequence, bounds its energy by (9.11), extracts a psh subnet by compactness, and uses Proposition 9.19 and energy upper semicontinuity to prove attainment. Corollary 9.12 identifies its maximizer with a genuine MA solution using the energy derivative. Proposition 12.7 gives injectivity modulo constants; Theorem 12.9 concludes existence and uniqueness. Restoring raw mass multiplies the measure by \(V\), exactly as in (1)--(2). There is no assumed torus action or skeletal support at any step. Thus T is an independent general-dimension arbitrary-support result, not evidence that D or a densely valued field automatically inherits it.

\section*{8. Reference independence and the exact partial theorem}

In D, if another semipositive model reference differs by a bounded model potential \(h\), then \(\theta'=\theta+dd^ch\) and the same metric is represented by \(u'=u-h\). Model curvature identities, the energy cocycle, and decreasing continuity give
\[
 M_{\theta'}(u-h)=M_\theta(u),\qquad
 E_{\theta'}(u-h)=E_\theta(u)-E_\theta(h).
 \tag{18}
\]
To check the unbounded passage, model majorants \(u_\alpha\downarrow u\) transform to \(u_\alpha-h\downarrow u-h\); both energy and MA are decreasing-continuous. The test cones correspond by the same translation. Consequently
\[
 Q_{\theta'}(\nu)=Q_\theta(\nu)+\int h\,d\nu-E_\theta(h),
 \quad
 \mathscr E_{\theta'}(\mu)=\mathscr E_\theta(\mu)+\int h\,d\mu-VE_\theta(h).
 \tag{19}
\]
All added terms are finite. Hence finite energy and the raw equation are reference-independent in the stated class. The same check uses the intersection-pairing cocycle for model/test-configuration references in T. No singularity type of a big or nonample class is substituted for the ample reference.

The partial theorem is therefore: separately in D and T, the genuine operator on \(\mathcal E^1(\theta)\) is independent of decreasing model approximations, has raw mass \((L^n)\), obeys comparison, and has range exactly (1). Every measure in that range has a solution unique modulo constants. There is no support restriction, toric invariance, or dimension-one assumption. The D proof replaces the old algebraizability-based orthogonality input and the old support-based integral continuity by Sections 4 and 5, while retaining BFJ's geometric compactness and envelope inputs. The T result is directly source-established.

\section*{9. Concrete original obligations still unresolved}

The full P16.6 residual scope beyond these field endpoints is not settled by this document. Even retaining smoothness, geometrical integrality, projectivity and an ample line bundle, we have not proved the same arbitrary-support finite-energy theorem for every complete nontrivially valued, potentially nondiscrete field in the general-dimension noninvariant setting. No impossibility claim is made.

For such a field the supplied continuous operator and synthetic finite-energy measure theory leave the following indispensable tasks: establish compactness of the normalized global psh cone in a topology compatible with model evaluations; prove the needed envelope regularity or an actual substitute that realizes the limits as global psh potentials; construct the approximation-independent unbounded mixed operator with full raw mass and locality; justify finite-energy differentiation, not only test-function differentiation; prove comparison and uniqueness for that actual unbounded operator; and prove attainment or an equivalent realization theorem for every positive mass-\(V\), finite-dual-energy target. Some of these may follow from future inputs, but none is deleted by measure completeness alone.

The trivial-field source Theorem 12.8 explicitly makes potential completeness and MA surjectivity equivalent to an envelope property. Synthetic Theorem 2.26 proves completeness of measures over arbitrary fields, a different assertion. Its Remark 2.14 sets aside the needed potential theory. This is a precise logical obstruction to inferring the missing theorem, not a counterexample to it. Restricting targets to one skeleton, imposing toric invariance, switching to a different inverse measure, specializing to a curve, or using unproved base change would not discharge the original obligations.

No resolution, alteration, normalization or base-change preservation statement for an unbounded cone, energy, MA, or comparison is used. Singular spaces, positive residue characteristic, nonample classes, proper nonprojective spaces and nonalgebraizable spaces are the separate original atoms and are not claimed solved here.

\section*{10. Reproducibility, literal extracts, and evidence limits}

The original input is retained byte-identically as input.json and input-original.json, with decoded assignment-original.txt and task-original.json. Run \texttt{python3 audit\_sources.py} to verify the ten supplied PDF/text hashes, regenerate text directly from each actual PDF, and reproduce source-extracts/source-provenance.json with literal published proof blocks, versions, identities, application checks and limitations. The audit output in source-audit-output.txt records ten successful PDF re-extractions and all supplied hashes matching. The independently regenerated text is not substituted for the original cache identity.

Rendered actual-PDF pages retained in source-extracts/synthetic-page25.png, source-extracts/bfj-v1-page27.png and source-extracts/bj-v4-page122.png resolve the absolute values in synthetic (2.27), the full-mass truncation estimate, and the exact trivial-field theorem/proof respectively. The additional source-extracts/bfj-v1-page28.png verifies that the sign slip in Lemma 6.10 is present in the PDF itself; (7a)--(7b) supply the required bound. The provenance file records these inspections. Rerun page rendering and the two builds with \texttt{python3 build\_artifacts.py}. That script regenerates a standalone P16.6-paper.tex from this complete argument and executes pdflatex twice in separate directories. Each actual stdout/stderr diagnostic stream, TeX log, command, exit code and hash is retained. The final PDF is P16.6-paper.pdf. Environment/tool versions are recorded in environment.json. An initial missing-package declaration for the mathscr command caused two failed TeX attempts; their actual diagnostics and source snapshot are retained in build/failed-missing-mathscr. The final preamble declares the existing mathrsfs package; no software was installed.

status.json has exactly the keys id, status, exact\_result, proof\_file and concrete\_gaps. No heavy computation, supervised-compute job, network fetch, software installation, other model or subagent was used. These checks establish reproducibility and artifact/source identity, not mathematical correctness. The candidate flag remains false regardless of build success.

\section*{References with exact supplied versions}

[BFJ-PSH] S. Boucksom, C. Favre, M. Jonsson, Singular semipositive metrics in non-Archimedean geometry, arXiv:1201.0187v2.

[BFJ-MA-v1] S. Boucksom, C. Favre, M. Jonsson, Solution to a non-Archimedean Monge-Ampere equation, arXiv:1201.0188v1. Theorem A is not read as a statement of a later version.

[BE] S. Boucksom, D. Eriksson, Spaces of norms, determinant of cohomology and Fekete points in non-Archimedean geometry, arXiv:1805.01016v3.

[BJ-T] S. Boucksom, M. Jonsson, Global pluripotential theory over a trivially valued field, arXiv:1801.08229v4.

[BJ-S] S. Boucksom, M. Jonsson, Measures of finite energy in pluripotential theory: a synthetic approach, arXiv:2307.01697v1. Its Example 2.20 is the supplied source for the published-background BE21/BGM22 input.

[PS] L. Pille-Schneider, Global pluripotential theory on hybrid spaces, arXiv:2209.04879v2, especially Remark 4.4.

[GR] W. Gubler, J. Rabinoff, Classically psh and pluriharmonic functions on Berkovich spaces, arXiv:2506.13548v2.

[Thuillier] A. Thuillier, supplied thesis HAL tel-00010990v1.

[Berkovich] V. G. Berkovich, Etale cohomology for non-Archimedean analytic spaces, Publications mathematiques de l'IHES 78 (1993), 5--161, supplied Numdam PDF.

[CLD] A. Chambert-Loir, A. Ducros, Formes differentielles reelles et courants sur les espaces de Berkovich, arXiv:1204.6277v2, July 2025 revision.
\end{document}
